\documentclass[10pt,a4paper]{amsart}
\usepackage[margin=19mm]{geometry}
\usepackage{amsmath,amssymb,mathtools}
\usepackage[scr=rsfs,cal=euler]{mathalfa}
\usepackage{xfrac}
\usepackage[unicode,hidelinks,bookmarks=false]{hyperref}
\newcommand{\ie}{i.e.\ }
\newcommand{\cf}{cf.\ }
\newcommand{\resp}{respectively\ }
\newcommand{\Subsection}{\subsection}
\providecommand{\nobreakdash}{\nobreak\hbox{-}}
\setlength{\emergencystretch}{2em}
\multlinegap0pt
\usepackage{bm}
\usepackage{bbm}
\usepackage{dsfont}
\usepackage{xspace}
\usepackage{cancel}
\usepackage{mathtools}
\usepackage{amsthm}
\makeatletter
\@ifundefined{theorem}{\newtheorem{theorem}{Theorem}}{}
\@ifundefined{lemma}{\newtheorem{lemma}[theorem]{Lemma}}{}
\@ifundefined{proposition}{\newtheorem{proposition}[theorem]{Proposition}}{}
\makeatother
\newcommand{\F}{\mathbb{F}}
\newcommand{\Q}{\mathbb{Q}}
\newcommand{\R}{\mathbb{R}}
\newcommand{\card}[1]{\lvert #1\rvert}
\newcommand{\dist}{\operatorname{d}}
\newcommand{\inner}[2]{\left\langle #1,#2\right\rangle}
\newcommand{\one}{\mathbf{1}}
\newcommand{\rank}{\operatorname{rank}}
\newcommand{\supp}{\operatorname{supp}}
\title[A Sharp Local-Question Threshold for GHZ-Equatorial Complete]{A Sharp Local-Question Threshold for GHZ-Equatorial Completeness in Four-Player XOR Games}
\author{Ziao Tang, Chengkai Zhu, Ge Bai, Xin Wang,, Ranyiliu Chen}
\date{Source: arXiv:2608.11139v1 (CC BY 4.0)}
\begin{document}
\maketitle

\noindent\emph{Source and licence.} A Sharp Local-Question Threshold for GHZ-Equatorial Completeness in Four-Player XOR Games. Version arXiv:2608.11139v1 (CC BY 4.0), distributed under the Creative Commons Attribution 4.0 International licence, as recorded in the arXiv licence metadata for this exact version and shown on the abstract page https://arxiv.org/abs/2608.11139v1. Full rights notice: licenses/arxiv-2608.11139-CC-BY-4.0.txt.\par\smallskip

\noindent\textit{Editorial scope.} Counted unit: the Proposition whose source label is \texttt{prop:ten-cover} in the article (source0.tex lines 1463--1466, source label \texttt{prop:ten-cover}), delivered together with the article's own complete original proof of it (source0.tex lines 1468--1557, one \texttt{proof} environment of 90 lines). No mathematics is written, repaired or re-derived: statement and proof are the article's own text line for line. Required context retained uncounted: lem:no-singleton-block (lines 591--596); prop:rank-cap (lines 605--617). Every definition, display and label of those lines is retained unchanged. Omitted from this package and not redistributed: 1--590, 597--604, 618--1462, 1467--1467, 1558--1980. Nothing inside a retained excerpt is omitted or altered: every retained body is delivered exactly as the article's lines are written, so no omission receipt of any kind accompanies this package. Citations: the retained excerpts carry no citation command at all, so no bibliography entry of the article is transcribed and no reference is redistributed. Reference adaptation: none was needed and none was made; every reference inside the delivered unit resolves inside it, so no unresolved reference or citation is produced. Printed slips kept literal: no printed word, symbol, hypothesis or number of the counted statement or of its proof was altered. Layout and class: the article's own class is \texttt{article} with packages including fontenc, lmodern, amsmath, amssymb, amsthm, mathtools, bm, fullpage, booktabs, longtable, array, tikz, natbib, microtype; the wrapper is the lane's pinned \texttt{amsart} editorial wrapper at 10pt on A4, which restates the theorem-like environments the retained text needs and adds the article's own macro definitions that the retained text needs (\texttt{\textbackslash F}, \texttt{\textbackslash Q}, \texttt{\textbackslash R}, \texttt{\textbackslash card}, \texttt{\textbackslash dist}, \texttt{\textbackslash inner}, \texttt{\textbackslash one}, \texttt{\textbackslash rank}, \texttt{\textbackslash supp}), with extra wrapper packages (bm, bbm, dsfont, xspace, cancel, mathtools). The delivered numbering is the unit's own. Assets: no retained span contains a figure environment, a graphics-inclusion command, a PDF-page inclusion, a table or a TikZ picture; no image file is copied into this package, the package declares no asset dependency and no image file is redistributed with it. Licence: version arXiv:2608.11139v1 of this article is distributed under the Creative Commons Attribution 4.0 International licence, as recorded in the arXiv licence metadata for this exact version and shown on the abstract page https://arxiv.org/abs/2608.11139v1. A full rights notice is delivered at licenses/arxiv-2608.11139-CC-BY-4.0.txt. Characters: every retained excerpt is pure ASCII, so the delivered engine, XeLaTeX with its default Latin Modern text font, reports no missing character.
\begin{lemma}
\label{lem:no-singleton-block}
Let \(c\in K_E\) be a circuit with support \(S=\supp c\).  Every
player-question block that is active on \(S\) contains at least two clauses
of \(S\).
\end{lemma}

\begin{proposition}\label{prop:rank-cap}
If \(c\) is a primitive circuit on \(S\), then, for the restriction \(A_S\)
of \(A\) to the rows indexed by \(S\),
\[
  \rank_{\Q}A_S=\card S-1,
  \qquad
  \card S\leq
  2+\sum_{\alpha=1}^{4}
  \bigl(\card{\{q_e^{(\alpha)}:e\in S\}}-1\bigr)\leq10.
\]
At ten clauses all pages are ternary; at nine clauses the question-count
profile is, up to permutation, \((3,3,3,2)\) or \((3,3,3,3)\).
\end{proposition}

\begin{proposition}\label{prop:ten-cover}
Ten distinct radius-one balls covering \([3]^4\) cannot have centers that
support a primitive circuit.
\end{proposition}

\begin{proof}
Fix an identification of the three question labels on each page with
\(\F_3\).  This choice does not affect Hamming distances.  Let \(S\) be the
set of ten centers and consider the ternary \([4,2,3]_3\)
Hamming code
\[
  C=\{(a,b,a+b,a+2b):a,b\in\F_3\}.
\]
Any two zero coordinates in a codeword force \(a=b=0\); hence every nonzero
codeword has weight at least three.  Since \((1,0,1,1)\in C\), the displayed
code has minimum distance three.  Its nine cosets partition \(\F_3^4\) into
codes of size nine.

Every radius-one ball meets each coset exactly once.  Indeed, the zero error
and the eight nonzero errors supported on one coordinate represent distinct
cosets of \(C\), because the difference of two such errors has weight at most
two.  Since \(\F_3^4/C\) has nine elements, these errors form a complete set
of coset representatives; translation by a ball center proves the claim.

Ten balls contribute ten incidences to each nine-point coset.  If they cover
the cube, every point has positive multiplicity, so each coset contains
exactly one double point and eight single points.  Across the nine cosets
there are exactly nine double points and no point of higher multiplicity.  If
\(n_d\) denotes the number of unordered pairs of centers at distance \(d\),
two radius-one balls meet in three points at center distance one, in two
points at distance two, and not at all at larger distances.  Counting the
nine double points gives
\begin{equation}\label{eq:intersection-app}
  3n_1+2n_2=9.
\end{equation}

Suppose \(S\) supports a primitive circuit.  Every page is ternary and,
by Lemma~\ref{lem:no-singleton-block}, every question block has size at least
two.  Among partitions of ten into three parts of size at least two,
\((4,3,3)\) minimizes the number of within-block pairs, giving
\[
  \binom42+\binom32+\binom32=12.
\]
Summed over four pages, the number of coordinate agreements satisfies
\[
  3n_1+2n_2+n_3
  \geq48.
\]
Together with \(n_1+n_2+n_3+n_4=45\), this implies
\[
  2n_1+n_2-n_4
  \geq3.
\]

With \(e_0,e_1,e_2\) denoting the standard basis of \(\R^3\), define
\[
  \Phi(q)=
  \left(
  \tfrac1{\sqrt3},
  e_{q_1}-\tfrac13\one,
  e_{q_2}-\tfrac13\one,
  e_{q_3}-\tfrac13\one,
  e_{q_4}-\tfrac13\one
  \right).
\]
The Gram matrix of the ten feature vectors satisfies
\[
  \mathcal G_{ij}
  =
  \inner{\Phi(q_i)}{\Phi(q_j)}
  =
  3-\dist(q_i,q_j).
\]
Centered and raw question indicators have the same column span after the
constant column is included.  Proposition~\ref{prop:rank-cap} therefore gives
\(\rank \mathcal G=\rank A_S=9\), and \(\mathcal G\) and
\(A_S^{\mathsf T}\) have the same
kernel on \(\R^S\).

The off-diagonal nonzero graph of \(\mathcal G\) joins pairs at distances one,
two, and four.  It must be connected.  Otherwise \(\mathcal G\) is block
diagonal, and
restricting the full-support circuit vector to each connected component gives
a nonzero kernel vector on that component.  Two components would yield two
independent kernel vectors, contrary to the one-dimensional circuit kernel.
A connected graph on ten vertices has at least nine edges, whereas
Eq.~\eqref{eq:intersection-app} gives
\[
  n_1+n_2+n_4
  =
  9-(2n_1+n_2-n_4)
  \leq6.
\]
This contradiction proves the claim.
\end{proof}

\end{document}
